% ---------------------------------------------------------------- experience
\section{Experience}\label{sec:exp}

% ............................................................ VisionRise
\entry[exp:visionrise]{VisionRise Technosoft Pvt.\ Ltd.}{Apr 2026 -- Oct 2026 · Mumbai}
\dek{Software Engineer Intern. Built the company's white-label learning platform and ran its production server.}

\begin{facts}
\fact{Company}{Edtech and software company behind the Rise Academy / Rise Institute and eGuru brands
  (\href{https://risetechnosoft.com}{risetechnosoft.com})}
\fact{Ownership}{Created and wrote nearly all of the LMS v2 codebase and the eGuru fork; contributed to
  the earlier, team-built LMS v1}
\fact{Stack}{TypeScript, Node.js, Express, PostgreSQL, Drizzle ORM, Zod, JWT, React Native, Expo,
  FFmpeg, Google Cloud Storage, Gemini, Zoom, Docker, Coolify, Ubuntu 24.04 LTS}
\fact{Live}{\href{https://play.google.com/store/apps/details?id=tech.riseinstitute.app}{Rise Academy}
  and \href{https://play.google.com/store/apps/details?id=app.eguru.lms}{eGuru} on Google Play
  \,·\, \href{https://riseinstitute.tech}{riseinstitute.tech} \,·\, \href{https://eguru.app}{eguru.app}}
\end{facts}

\runin{Context} LMS v2 is a white-label learning platform: one backend and one app codebase that
ship under more than one brand, covering courses, live classes, assessments and paid video for
students on Android, iOS and the web. Kartik built it from the database schema to the store
listing, then forked it into eGuru, a second branded app.

\runin{Backend architecture} A Node.js/Express API in TypeScript (\tld40,000 lines), organised into
20 domain modules that together expose 276 REST routes: admin, AI, analytics, assignments, audit,
auth, calendar, certificates, courses, enrollments, lessons, live classes, media, modules,
notifications, payments, quizzes, reports, settings and users. Inputs are validated with Zod at the
boundary, data lives in PostgreSQL with schema changes tracked as Drizzle migrations, and the
repository carries 25 test files.

\runin{One app, three platforms} The client is a single React Native/Expo codebase (\tld45,000
lines) that builds for Android, iOS and the web, with push notifications and over-the-air updates
through Expo, so fixes reach installed apps without waiting on a store release.

\runin{Protecting paid content} Paid video is the asset most worth protecting. Uploads go through an
FFmpeg transcoding pipeline that produces AES-128-encrypted HLS. Playback needs a per-session token
that the API validates on every segment request, and media is fetched through short-lived signed
Google Cloud Storage URLs, so a copied link expires quickly and does not work outside the session
that requested it. PDFs stream inside the app with no download surface, and course certificates
can be verified through a SHA-256 hash.

\begin{center}
\vspace{2pt}
\begin{tikzpicture}[
  font=\sffamily\footnotesize,
  box/.style={draw=hair, line width=0.5pt, rounded corners=2pt, fill=wash, text=ink,
              align=center, inner xsep=5pt, inner ysep=4pt, minimum width=36mm, minimum height=10mm},
  arrow/.style={-{Stealth[length=4pt,width=3.5pt]}, draw=muted, line width=0.55pt},
  tag/.style={font=\sffamily\scriptsize, text=muted}]
  \node[box] (up)  at (0,0)      {Course video\\uploaded by staff};
  \node[box] (ff)  at (5.1,0)    {FFmpeg transcode\\to AES-128 HLS};
  \node[box] (gcs) at (10.2,0)   {Google Cloud Storage\\encrypted segments};
  \node[box] (app) at (0,-1.75)  {Student app\\React Native / Expo};
  \node[box] (api) at (5.1,-1.75){API checks per-session\\token on every segment};
  \node[box] (url) at (10.2,-1.75){Short-lived\\signed URL};
  \draw[arrow] (up) -- (ff);
  \draw[arrow] (ff) -- (gcs);
  \draw[arrow] (app) -- node[tag, above]{request} (api);
  \draw[arrow] (api) -- node[tag, above]{issue} (url);
  \draw[arrow] (url) -- (gcs);
\end{tikzpicture}\par
\vspace{3pt}
{\footnotesize\itshape\color{muted}Paid-video path: content is encrypted when it is transcoded and is reachable only with a session token and an expiring URL.}
\end{center}

\runin{Authentication and access} Students sign in with Google or with a PIN tied to a phone number
verified over WhatsApp. Sessions use JWT access and refresh tokens, role-based access separates
admins, instructors and students, and admin actions are written to an audit log.

\runin{Features and integrations} Kartik built Live Classes v1 on Zoom: classes are orchestrated from
the platform, Zoom's webhooks are signature-verified before they are trusted, attendance and
recordings are captured, and Gemini then generates a quiz from the class content. Around this core
Kartik shipped courses, quizzes and assignments, UPI payment verification, push and email
notifications, and bulk enrollment.

\runin{Production server} Kartik provisioned and hardened the company's primary production server on
Ubuntu 24.04 LTS with a standard baseline: SSH key-only access, a firewall, unattended security
upgrades and fail2ban. The apps and PostgreSQL run behind Coolify, which handles CI/CD and
deployments; container and application logs are collected through Coolify and Docker; and automated
backups go to Google Cloud Storage.

\runin{IT administration} Kartik administered the company's GitHub organisation (15 repositories,
about 5--10 members) and its Google Workspace directory of 35+ accounts, covering accounts, access
and permissions.

\runin{Other work} Contributed to the earlier, team-built LMS v1: CI/CD, timezone-safe class
scheduling with the PostgreSQL session pinned to IST, course-invite and access-request flows, and
Zoom fixes. Built a government procurement-management prototype (FastAPI and React) for a client
demo, and the eguru.app marketing site as a Next.js static export.

\runin{Outcome} Two apps live on Google Play, Rise Academy and eGuru, built from one codebase and one
backend, and deployed through a pipeline onto a server that Kartik set up and maintained.

% ............................................................ R.U.D.R.A
\entry[exp:rudra]{R.U.D.R.A Cybersecurity Pvt.\ Ltd.}{May 2025 -- Jul 2025 · Mumbai}
\dek{Security Engineer Intern for three months at a Mumbai-based VAPT and offensive-security firm: client testing, internal tooling and purple-team exercises, using BBOT, Semgrep, Neo4j, Nuclei and Sliver C2.}

\runin{Attack-surface monitoring} Kartik automated recurring external attack-surface monitoring with
BBOT, cutting manual reconnaissance and saving \tld\$20,000 a year.

\runin{Static analysis at scale} Kartik single-handedly built a custom Semgrep (SAST) module with
Neo4j-based graph visualisation, for analysing large codebases.

\runin{Banking client} Manual code review of a banking client's web portal found two high-severity
vulnerabilities: a local file inclusion, and broken access control that exposed restricted
functionality to unauthorised users. Both were reported to the client, and Kartik then automated
their detection with custom Nuclei templates.

\runin{Purple teaming} Deployed and operated the Sliver C2 framework for purple-team exercises that
emulated real-world intrusions.

\runin{Outcome} Two high-severity findings for a client, a recurring recon task automated with a
\tld\$20,000-a-year saving, and an in-house SAST module.

% ............................................................ Freelance
\entry[exp:freelance]{Freelance}{May 2025 -- Oct 2026}
\dek{Self-employed Backend \& Security Engineer. Two separate engagements, not continuous; clients are not named.}

\minientry{Incident response and hardening}{May -- Jun 2025 · about one month}
A freelance client ran a web platform with 16,000+ users. Kartik hardened its server
infrastructure, and when the platform was hit by ransomware, led the incident response: contained
the attack and restored service through rapid backup and recovery.

\minientry[proj:sports]{sports-live: tournament auction and scoring platform}{Aug 2026 -- Oct 2026}
Built for a client and used at ELEV8, a pickleball tournament in which \tld200 players were
auctioned into 8 teams. The platform covers player registration and verification, a purse-based
live player auction with a board for the projector, fixtures, live scoring and standings. Venue
Wi-Fi is unreliable, so writes go into an offline queue in IndexedDB and are sent when the
connection returns.

\begin{facts}
\fact{Ownership}{Solo}
\fact{Stack}{Next.js 16, Supabase (PostgreSQL with row-level security, Realtime, Storage),
  server-side image processing, IndexedDB, Vitest}
\end{facts}
